\section{Oxide}

\subsection{Phân loại oxide}

Oxide là hợp chất của oxygen với một nguyên tố khác. Ở mức THCS, thường gặp hai nhóm chính:

\begin{itemize}
\item \textbf{oxide base}: thường là oxide của kim loại, chẳng hạn \ce{CaO}, \ce{MgO}, \ce{CuO}, \ce{Fe2O3};
\item \textbf{oxide acid}: thường là oxide của phi kim, chẳng hạn \ce{CO2}, \ce{SO2}.
\end{itemize}

\begin{ghinho}{Mối liên hệ cần nhớ}
Oxide base thường phản ứng với acid tạo muối và nước. Một số oxide base phản ứng với nước tạo base.

Oxide acid thường phản ứng với base tạo muối và nước. Một số oxide acid phản ứng với nước tạo acid.
\end{ghinho}

Ví dụ:
\[
\ce{CuO + 2HCl -> CuCl2 + H2O},
\]
\[
\ce{CaO + H2O -> Ca(OH)2},
\]
\[
\ce{CO2 + Ca(OH)2 -> CaCO3 v + H2O},
\]
\[
\ce{SO2 + H2O -> H2SO3}.
\]

\subsection{Dự đoán phản ứng và hiện tượng}

Khi đề mô tả một oxide phản ứng với dung dịch, hãy làm theo thứ tự:
\begin{enumerate}
\item xác định oxide thuộc loại nào;
\item xác định chất còn lại là acid hay base;
\item dự đoán sản phẩm;
\item kiểm tra hiện tượng có phù hợp hay không.
\end{enumerate}

\begin{vidu}
Dẫn khí \ce{CO2} vào nước vôi trong. Nêu hiện tượng khi lượng \ce{CO2} vừa đủ và viết phương trình hóa học.
\end{vidu}
\begin{loigiai}
\ce{CO2} là oxide acid, phản ứng với \ce{Ca(OH)2} tạo \ce{CaCO3} không tan:
\[
\ce{CO2 + Ca(OH)2 -> CaCO3 v + H2O}.
\]
Dung dịch xuất hiện kết tủa trắng.

\textbf{Vậy,} nước vôi trong trở nên đục do tạo kết tủa \ce{CaCO3}.
\end{loigiai}

\subsection{Một số phản ứng thường gặp}

Khi \ce{CO2} được dẫn vào dung dịch \ce{Ca(OH)2} với lượng dư, kết tủa \ce{CaCO3} có thể tan do tạo bicarbonate tan:
\[
\ce{CaCO3 + CO2 + H2O -> Ca(HCO3)2}.
\]
Phản ứng này giúp giải thích một số hiện tượng liên quan đến nước cứng và sự hòa tan đá vôi trong tự nhiên.

\begin{tuluyen}
Cho các oxide: \ce{CaO}, \ce{CO2}, \ce{CuO}, \ce{SO2}. Phân loại thành oxide acid và oxide base.
\dapso{Oxide base: \ce{CaO}, \ce{CuO}; oxide acid: \ce{CO2}, \ce{SO2}.}
\begin{loigiaiBT}
\ce{CaO}, \ce{CuO} là oxide của kim loại và thuộc nhóm oxide base. \ce{CO2}, \ce{SO2} là oxide của phi kim và thuộc nhóm oxide acid.

\textbf{Vậy,} oxide base là \ce{CaO}, \ce{CuO}; oxide acid là \ce{CO2}, \ce{SO2}.
\end{loigiaiBT}
\end{tuluyen}
